\subsection{Literature}  \label{sec:lit_review}

This paper connects three lines of research: infrastructure capitalization, local public finance, and the measurement of infrastructure quality.
The common thread is that households do not buy homes separate from the local public-service bundle. Along with a house, they buy access to a jurisdiction's taxes, services, and expected public infrastructure maintenance choices. Consequently, road-maintenance renewal elections are both fiscal decisions and housing-market signals.

{\bf Infrastructure and house values.} A large literature studies how transportation infrastructure affects local housing markets. Much of this work focuses on new infrastructure or major upgrades. For example, \cite{gonzalez2016paving} studies randomized street paving in Mexico and finds large increases in property values, while \cite{theisen2021road} and related work estimate capitalization effects from new or improved transportation access \citep{dube2018bus}. These studies are important, but they address a different policy margin than the one facing many U.S. local governments. In mature road networks, the central question is often not where to build a new road but whether to preserve existing roads before deferred maintenance becomes costly and visible.

The maintenance margin is less studied, especially for local roads in developed economies. \cite{gertler2024road} examines highway maintenance in Indonesia and finds that better road quality affects local economic activity and housing prices. The classic capitalization study of \cite{edel1974taxes} examines local public spending and house values in the Boston area but predates modern quasi-experimental designs and high-frequency administrative housing data. Our contribution is to study local road maintenance in a setting where existing revenue streams are periodically renewed or cut by voters and where close elections provide quasi-experimental variation.

{\bf Local public finance and ballot-box decisions.} The local public finance tradition emphasizes that property values capitalize local taxes and public services \citep{Oates1969, Oates1972}. Recent work on local tax relief has documented the immediate capitalization effects of property tax cuts \citep{moulton2018benefits}, but it has not focused on the long-term effects of cutting dedicated maintenance funding. Renewal levies make this trade-off explicit. If voters renew a levy, they continue paying a dedicated tax and preserve maintenance funding. If they reject it, they receive near-term tax relief but reduce the local government's capacity to maintain a public asset. This setting is related to the use of referendum discontinuities in public finance, including \cite{cellini2010value}, but differs in an important way: we study renewals for existing maintenance rather than new capital spending. Renewal timing is set by the expiration of prior levies, which helps separate the maintenance-funding shock from the endogenous timing of new investment proposals.

{\bf Measuring road quality.} A key challenge to evaluating infrastructure policy is measurement. Engineering measures such as the International Roughness Index and Pavement Condition Index can be costly, irregular, and unavailable for many local roads. \cite{currier2023} shows how vertical-acceleration signals can be used to construct road-roughness measures in places with sufficient platform coverage, while \cite{wong2017local} relies on standardized field surveys to measure road quality in rural China. Our approach is different and complements these existing methods. We fine-tune a vision model on labeled road samples to construct consistent road-quality measures for Ohio jurisdictions. This approach allows us to build a panel of road images to observe whether lost maintenance revenue is followed by road deterioration.



